\section{Results}
\label{sec:results}

This section summarizes the currently completed 2015, 2016 10\%, and 2015+2016
combination outputs. The quoted figures are still pre-rerun products in the sense
described earlier: they were generated before the final leakage-retaining template
update and before the optional radiative-fraction penalty was introduced. They are
therefore best read as the present physics organization of the result rather than as
the last word on the numerical limits. Even so, they already show the structure of the
analysis clearly and indicate why the global-GP approach is promising for a combined
prompt $A^\prime$ search.

\subsection{2015 analysis}
\label{sec:results-2015}

The 2015 analysis uses the fully unblinded engineering-run spectrum over
\SIrange{20}{130}{MeV}. Relative to the published local-window search, the present GPR
study is being validated over a substantially wider mass interval while preserving a
smooth, well-behaved upper-limit curve. This is an important methodological point: the
value of the GP approach is not merely that it fits the local background, but that it
does so in a way that remains stable when the search window is widened and the local
fit boundaries of the earlier polynomial approach are removed.

\begin{figure}[H]
\centering
\begin{subfigure}[t]{0.48\linewidth}
  \centering
  \graphicorplaceholder{\linewidth}{upper_limit_figs/2015/ul_bands_signal_yield_obsexp.png}
  \caption{Signal-yield upper limits with expected bands.}
\end{subfigure}
\hfill
\begin{subfigure}[t]{0.48\linewidth}
  \centering
  \graphicorplaceholder{\linewidth}{upper_limit_figs/2015/ul_bands_eps2_obsexp.png}
  \caption{$\eps^2$ upper limits with expected bands.}
\end{subfigure}
\caption{2015 upper limits for the current 95\% CL workflow. The left panel is the
limit in total signal-yield space and the right panel is the corresponding mapping to
$\eps^2$. These curves should be interpreted as pre-rerun baselines, but they already
demonstrate the smooth scan structure obtained with the global-GP background model.}
\label{fig:results-2015-ul}
\end{figure}

\begin{figure}[H]
\centering
\begin{subfigure}[t]{0.48\linewidth}
  \centering
  \graphicorplaceholder{\linewidth}{significance_figs/2015/ul_pvalues_components_local_global_refs.png}
  \caption{Toy-tail significance diagnostics: $p_{\rm strong}$, $p_{\rm weak}$, and
  $p_{\rm two}$.}
\end{subfigure}
\hfill
\begin{subfigure}[t]{0.48\linewidth}
  \centering
  \graphicorplaceholder{\linewidth}{significance_figs/2015/p0_analytic_local_global.png}
  \caption{Analytic local and provisional global $p_0$ curves.}
\end{subfigure}
\caption{2015 significance summary. The left panel reports where the observed limit
falls relative to the toy-limit ensemble, while the right panel shows the analytic
profile-likelihood significance curves. The global component remains provisional until
the final look-elsewhere treatment is rerun.}
\label{fig:results-2015-significance}
\end{figure}

\subsection{2016 10\% analysis}
\label{sec:results-2016}

The 2016 result shown here corresponds to the 10\% development subset used for the
present statistical validation program. The mass resolution, radiative-fraction
approximation, and search interval all differ from the 2015 case, but the statistical
machinery is identical. In particular, the significance curves now carry the full set
of analytic and toy-tail diagnostics discussed in Section~\ref{sec:methodology}, which
makes the interpretation of upward and downward fluctuations considerably more
transparent than in the earliest draft studies.

\begin{figure}[H]
\centering
\begin{subfigure}[t]{0.48\linewidth}
  \centering
  \graphicorplaceholder{\linewidth}{upper_limit_figs/2016_10pct/ul_bands_signal_yield_obsexp.png}
  \caption{Signal-yield upper limits with expected bands.}
\end{subfigure}
\hfill
\begin{subfigure}[t]{0.48\linewidth}
  \centering
  \graphicorplaceholder{\linewidth}{upper_limit_figs/2016_10pct/ul_bands_eps2_obsexp.png}
  \caption{$\eps^2$ upper limits with expected bands.}
\end{subfigure}
\caption{2016 10\% upper limits at 95\% CL. These are development-stage results, but
they already show the expected smooth behavior of the global-GP procedure across the
expanded 2016 mass interval.}
\label{fig:results-2016-ul}
\end{figure}

\begin{figure}[H]
\centering
\begin{subfigure}[t]{0.48\linewidth}
  \centering
  \graphicorplaceholder{\linewidth}{significance_figs/2016_10pct/ul_pvalues_components_local_global_refs.png}
  \caption{Toy-tail significance diagnostics.}
\end{subfigure}
\hfill
\begin{subfigure}[t]{0.48\linewidth}
  \centering
  \graphicorplaceholder{\linewidth}{significance_figs/2016_10pct/p0_analytic_local_global.png}
  \caption{Analytic local and provisional global $p_0$ curves.}
\end{subfigure}
\caption{2016 10\% significance summary. The separation between toy-tail
diagnostics and analytic profile-likelihood significance is important: the former test
how unusual the observed \emph{limit} is relative to the toy ensemble, while the latter
test how signal-like the observed spectrum is under the local extraction model.}
\label{fig:results-2016-significance}
\end{figure}

\subsection{2015+2016 10\% combined upper limit on \texorpdfstring{$\eps^2$}{epsilon2}}
\label{sec:results-combined-1516}

The first shared-$\eps^2$ combination joins the fully unblinded 2015 spectrum with the
2016 10\% development sample. This is already a meaningful milestone for the analysis
because it shows that the GPR framework can combine statistically independent data sets
with different resolutions, search ranges, and radiative normalizations inside a single
likelihood rather than by comparing precomputed limit curves after the fact.

The combined parameter of interest is $\eps^2$, not the individual signal yield of a
single data set. For that reason the combined result is quoted directly as an
$\eps^2$ limit. This is also the natural place to compare the current GPR reach to the
published HPS prompt exclusions once the final rerun is complete and the comparison
overlay has been frozen.

\begin{figure}[H]
\centering
\graphicorplaceholder{0.86\linewidth}{upper_limit_figs/2015_2016_combined/ul_bands_eps2_obsexp.png}
\caption{Combined 2015+2016 10\% upper limit on $\eps^2$ at 95\% CL. This is the first
result in the note that is intrinsically a shared-coupling measurement rather than a
single-spectrum reinterpretation. The present curve is still a pre-rerun baseline, but
it already demonstrates the expected gain from combining statistically independent
channels in a common coupling parameter.}
\label{fig:results-combined-1516-ul}
\end{figure}

\begin{figure}[H]
\centering
\begin{subfigure}[t]{0.48\linewidth}
  \centering
  \graphicorplaceholder{\linewidth}{significance_figs/2015_2016_combined/ul_pvalues_components_local_global_refs.png}
  \caption{Combined toy-tail significance diagnostics.}
\end{subfigure}
\hfill
\begin{subfigure}[t]{0.48\linewidth}
  \centering
  \graphicorplaceholder{\linewidth}{significance_figs/2015_2016_combined/p0_analytic_local_global.png}
  \caption{Combined analytic local and provisional global $p_0$ curves.}
\end{subfigure}
\caption{Combined significance summary for the 2015+2016 study. These plots are the
appropriate place to track how the shared-$\eps^2$ procedure behaves statistically
before the final global-significance evaluation is frozen.}
\label{fig:results-combined-1516-significance}
\end{figure}

\begin{figure}[p]
\centering
\graphicorplaceholder{0.96\linewidth}{combined_search_figs/extract_display_combined_m040MeV_z7p0.png}
\caption{Illustrative combined extraction display at \SI{40}{MeV}. The annotation box
summarizes the injected common coupling, the extracted coupling, the observed combined
limit, and the fitted per-dataset yields. The upper panels show how the two input data
sets respond individually, while the lower panel shows the combined Gaussian signal
representation. This figure is intentionally kept large because the interplay between
the dataset-level fits and the combined interpretation is the central visual message of
the shared-likelihood method.}
\label{fig:results-combined-1516-display}
\end{figure}

\analysisplaceholder{\textbf{Published-comparison placeholder.} Once the final rerun is
complete, insert an explicit comparison of the combined GPR limit with the published
2015 and 2016 prompt exclusions. That comparison will be especially informative because
the present GP-based analysis is already being validated over a wider search interval
than the earlier local-window searches.}

\subsection{2021 0.03\% analysis}
\label{sec:results-2021-003}

\analysisplaceholder{\textbf{Placeholder section.} Insert the 2021 0.03\% validation
upper-limit and significance products here once the corresponding summary suite is
finalized. This section should mirror the structure of the completed 2015 and 2016
subsections.}

\begin{figure}[H]
\centering
\begin{subfigure}[t]{0.48\linewidth}
  \centering
  \graphicorplaceholder{\linewidth}{upper_limit_figs/2021_0pt03pct/ul_bands_signal_yield_obsexp.png}
  \caption{Signal-yield upper limits.}
\end{subfigure}
\hfill
\begin{subfigure}[t]{0.48\linewidth}
  \centering
  \graphicorplaceholder{\linewidth}{upper_limit_figs/2021_0pt03pct/ul_bands_eps2_obsexp.png}
  \caption{$\eps^2$ upper limits.}
\end{subfigure}
\caption{Planned location for the 2021 0.03\% upper-limit summary.}
\label{fig:results-2021-003-placeholder}
\end{figure}

\subsection{2021 1\% analysis}
\label{sec:results-2021-1pct}

\analysisplaceholder{\textbf{Placeholder section.} This section is reserved for the
2021 1\% individual analysis and should ultimately include both the standalone result
and the way that subset enters the later shared-$\eps^2$ combinations.}

\begin{figure}[H]
\centering
\begin{subfigure}[t]{0.48\linewidth}
  \centering
  \graphicorplaceholder{\linewidth}{upper_limit_figs/2021_1pct/ul_bands_signal_yield_obsexp.png}
  \caption{Signal-yield upper limits.}
\end{subfigure}
\hfill
\begin{subfigure}[t]{0.48\linewidth}
  \centering
  \graphicorplaceholder{\linewidth}{upper_limit_figs/2021_1pct/ul_bands_eps2_obsexp.png}
  \caption{$\eps^2$ upper limits.}
\end{subfigure}
\caption{Planned location for the 2021 1\% upper-limit summary.}
\label{fig:results-2021-1pct-placeholder}
\end{figure}

\subsection{Combined limit on \texorpdfstring{$\eps^2$}{epsilon2} using 2015+2016 10\%+2021 1\%}
\label{sec:results-combined-151621-1pct}

This first three-sample study combines the full 2015 data set, the 2016 10\%
development subset, and the 2021 1\% subset using the most recent
leakage-retaining signal model. It should therefore be regarded as the current
post-leakage baseline for the shared-$\eps^2$ workflow. At the same time, it is
\emph{not} yet the final physics result: the radiative-fraction penalty has not
been applied, and the significance bookkeeping described below still needs one
more frozen rerun before any discovery-style statement is appropriate.

The observed combined $\eps^2$ limit in Figure~\ref{fig:results-combined-151621-1pct-ul}
is broadly well behaved and remains close to the expected band over most of the
scan. That general behavior is what one would expect from a statistically stable
combined analysis. The visible kinks in the curve should not be overinterpreted
as narrow resonant structure: this summary suite is labeled ``all'' because the
active data-set content changes across the scan. Only the approximate
\SIrange{35}{130}{MeV} interval is a genuine three-way overlap. Below that range,
the curve interpolates between 2015-only and 2015+2021 support, while above it
the result transitions to a 2016+2021 combination and finally to the 2021-only
tail near the upper edge of the plot. In that sense the limit plot looks
qualitatively reasonable.

The significance panels require more care. The toy-based quantities
$p_{\rm strong}$, $p_{\rm weak}$, and $p_{\rm two}$ in
Figure~\ref{fig:results-combined-151621-1pct-significance} are \emph{not}
discovery-significance estimators; they test how unusual the observed
\emph{upper limit} is relative to the background-only toy ensemble. By contrast,
the analytic $p_0$ curve is the local profile-likelihood bump-hunt statistic.
These two diagnostics need not agree numerically, and they should not be
extrapolated into one another. This is especially important here because the
toy-limit ensemble contains \num{10000} pseudoexperiments, so the smallest
resolvable one-sided tail probability is only $\mathcal{O}(10^{-4})$, i.e. about
\SI{3.7}{\sigma}; it cannot by itself validate or refute a
\SIrange{6}{7}{\sigma} excursion.

The striking feature of the present three-way summary is the narrow combined
analytic-significance spike near \SI{126}{MeV}. A combined local enhancement
larger than any individual input sample is qualitatively possible, because the
shared-$\eps^2$ fit pools statistically independent information with
mass-dependent weights. However, the auxiliary 2015, 2016, and 2021 local/global
panels in the same summary folder show only $\mathcal{O}(2$--$3)\sigma$ local
excursions and no comparably large global effect. For that reason, the current
combined peak is larger than one would expect from a naive extrapolation of the
toy-tail weak-$p$ values or from a simple quadrature estimate based on the
dataset-level panels. Under the statistical standards usually applied in
high-statistics HEP searches, that makes the feature interesting but not yet
claimable.

I verified that the \emph{form} of the global-significance conversion is the
standard one used elsewhere in this workflow: a one-sided local $p_0$ is mapped
to a global $p$-value with a \v{S}id\'{a}k trials-factor correction using an
effective $N_{\rm eff}$ derived from the scan range and the mass resolution.
That is an acceptable provisional look-elsewhere treatment. The caveat is that
the current summary-level significance products still need a frozen rerun before
their numerical values should be treated as final. In particular, the
individual-dataset global $p$-value overlays in this combined-summary folder
should be considered provisional even after restricting them to the
combined-support mass domain, and the present three-way combined significance
peak should be revisited after rerunning the corrected post-leakage summary with
the radiative-fraction penalties included. The correct conservative reading of
this study is therefore: the combined \emph{limit} behavior looks plausible and
encouraging, while the quoted global-significance values remain preliminary.

\begin{figure}[H]
\centering
\graphicorplaceholder{0.86\linewidth}{upper_limit_figs/2015_2016_2021_1pct_combined/ul_bands_eps2_obsexp.png}
\caption{Observed and expected 95\% CL upper limit on $\eps^2$ for the current
2015+2016 10\%+2021 1\% shared-coupling study using the leakage-retaining signal
update. The overall smoothness of the curve is consistent with a statistically
stable combination. Changes in slope are expected when the active set of input
data sets changes across the mass scan, so this figure should be read as an
``all available data'' summary rather than as a pure three-way-overlap result at
every mass point.}
\label{fig:results-combined-151621-1pct-ul}
\end{figure}

\begin{figure}[H]
\centering
\begin{subfigure}[t]{0.48\linewidth}
  \centering
  \graphicorplaceholder{\linewidth}{significance_figs/2015_2016_2021_1pct_combined/ul_pvalues_components_local_global_refs.png}
  \caption{Toy-tail diagnostics for the observed CL$_s$ limit.}
\end{subfigure}
\hfill
\begin{subfigure}[t]{0.48\linewidth}
  \centering
  \graphicorplaceholder{\linewidth}{significance_figs/2015_2016_2021_1pct_combined/p0_analytic_local_global.png}
  \caption{Analytic local and provisional global $p_0$ curves.}
\end{subfigure}
\caption{Current three-way combined significance summary. The left panel tests how
unusual the observed \emph{limit} is relative to background-only toy limits,
whereas the right panel shows the analytic bump-hunt statistic. The two should
not be identified with one another. The narrow high-significance structure near
\SI{126}{MeV} is the main feature requiring rerun-level scrutiny. The
look-elsewhere conversion itself follows the standard trial-factor logic, but
the numerical global $p$-values shown here, especially for the auxiliary
individual-dataset overlays in this study, should still be regarded as
preliminary until the post-leakage and radiative-penalty rerun is frozen.}
\label{fig:results-combined-151621-1pct-significance}
\end{figure}

\subsection{First-round unblinding: 2015+2016 10\%+2021 10\%}
\label{sec:results-first-unblinding}

\analysisplaceholder{\textbf{Placeholder section.} This section is reserved for the
first-round unblinding milestone based on the 2021 10\% workflow. It should include
the finalized local/global significance discussion only after the corrected
look-elsewhere rerun is available.}

\begin{figure}[H]
\centering
\begin{subfigure}[t]{0.48\linewidth}
  \centering
  \graphicorplaceholder{\linewidth}{upper_limit_figs/first_round_unblinding_2015_2016_2021_10pct/ul_bands_signal_yield_obsexp.png}
  \caption{Signal-yield upper limits.}
\end{subfigure}
\hfill
\begin{subfigure}[t]{0.48\linewidth}
  \centering
  \graphicorplaceholder{\linewidth}{upper_limit_figs/first_round_unblinding_2015_2016_2021_10pct/ul_bands_eps2_obsexp.png}
  \caption{$\eps^2$ upper limits.}
\end{subfigure}
\caption{Planned location for the first-round unblinding summary.}
\label{fig:results-first-unblinding-placeholder}
\end{figure}

\subsection{Individual 2016 analysis}
\label{sec:results-individual-2016-placeholder}

\analysisplaceholder{\textbf{Placeholder section.} Reserve this subsection for the
full or updated standalone 2016 analysis that supersedes the current 10\% development
subset.}

\subsection{Individual 2021 10\% analysis}
\label{sec:results-individual-2021-10pct-placeholder}

\analysisplaceholder{\textbf{Placeholder section.} Reserve this subsection for the
standalone 2021 10\% analysis products.}

\subsection{Fully unblinded 2021 data set analysis}
\label{sec:results-fully-unblinded-2021-placeholder}

\analysisplaceholder{\textbf{Placeholder section.} Reserve this subsection for the
fully unblinded 2021 data set once the final prompt selection, exposure, and global
significance inputs are frozen.}

\subsection{2015+2016+2021 complete combined analysis}
\label{sec:results-complete-combined-placeholder}

\analysisplaceholder{\textbf{Placeholder section.} Reserve this subsection for the
final combined HPS prompt-resonance result. The completed version should quote only the
validated analysis range, the corrected global-significance treatment, and the final
radiative-fraction inputs.}
